\section*{अध्याय \ref{ch:b1:structures} --- बीजीय संरचनाएँ}

\begin{solution}{exo:b1:structures:1}
\emph{स्थायित्व:} $x * y = 1 \iff x + y - xy = 1 \iff (1-x)(1-y) = 0$, जो
$x, y \neq 1$ के लिए असंभव है। वस्तुतः मुख्य सर्वसमिका यह है
\[
1 - x * y = (1 - x)(1 - y):
\]
\omterm{def:b1:logic:map}{प्रतिचित्रण} $\varphi(x) = 1 - x$ $(E, *)$ को $(\R^*, \times)$ पर भेजता है,
जहाँ $\varphi(x * y) = \varphi(x)\varphi(y)$ --- अर्थात् \omterm{def:b1:logic:inj}{एकैकी आच्छादक}
\omterm{def:b1:structures:morphism}{समाकारिता}। अब सभी अभिगृहीत पार चले जाते हैं: साहचर्यता और क्रमविनिमेयता
$\times$ की इन्हीं गुणों से निकलती हैं; तत्समक $\varphi^{-1}(1) = 0$ है
(जाँच: $x * 0 = x$); और $x$ का प्रतिलोम $\varphi^{-1}\bigl((1-x)^{-1}\bigr)
= 1 - \frac{1}{1-x} = \frac{x}{x - 1}$ है (जो $\neq 1$ है)। अतः $(E, *)$
\omterm{def:b1:structures:group}{आबेली समूह} है।
\end{solution}

\begin{solution}{exo:b1:structures:2}
\begin{enumerate}
  \item हाँ: धनात्मक संख्याओं का गुणनफल धनात्मक होता है, तत्समक $1$,
  प्रतिलोम $\frac 1x$, और साहचर्यता $\R^*$ से विरासत में।
  \item नहीं: स्थायी नहीं ($1 + 1 = 2 \notin \{-1,0,1\}$)।
  \item हाँ: यही मानक उदाहरण है।
  \item नहीं: स्थायी नहीं (विषम $+$ विषम $=$ सम), और कोई तत्समक भी नहीं ($0$
  सम है)।
\end{enumerate}
\end{solution}

\begin{solution}{exo:b1:structures:3}
$\mathrm{id}$ लिखिए, व्यत्यास $\tau_{12}, \tau_{13}, \tau_{23}$ (जो दोनों
नामित बिंदुओं की अदला-बदली करते हैं), और चक्र $c = (1\,2\,3)$ (अर्थात् $1
\mapsto 2 \mapsto 3 \mapsto 1$) तथा $c^2 = (1\,3\,2)$। $\sigma\rho$ की सारणी
(पंक्ति $\sigma$, स्तंभ $\rho$, पहले $\rho$ लगाइए):
\begin{center}
\renewcommand{\arraystretch}{1.2}
\begin{tabular}{c|cccccc}
$\sigma\backslash\rho$ & $\mathrm{id}$ & $c$ & $c^2$ & $\tau_{12}$ &
$\tau_{13}$ & $\tau_{23}$ \\ \hline $\mathrm{id}$ & $\mathrm{id}$ & $c$ &
$c^2$ & $\tau_{12}$ & $\tau_{13}$ & $\tau_{23}$ \\ $c$ & $c$ & $c^2$ &
$\mathrm{id}$ & $\tau_{13}$ & $\tau_{23}$ & $\tau_{12}$ \\ $c^2$ & $c^2$ &
$\mathrm{id}$ & $c$ & $\tau_{23}$ & $\tau_{12}$ & $\tau_{13}$ \\ $\tau_{12}$
& $\tau_{12}$ & $\tau_{23}$ & $\tau_{13}$ & $\mathrm{id}$ & $c^2$ & $c$ \\
$\tau_{13}$ & $\tau_{13}$ & $\tau_{12}$ & $\tau_{23}$ & $c$ & $\mathrm{id}$
& $c^2$ \\ $\tau_{23}$ & $\tau_{23}$ & $\tau_{13}$ & $\tau_{12}$ & $c^2$ &
$c$ & $\mathrm{id}$ \\
\end{tabular}
\end{center}
क्रम-विनिमेय न होने वाला युग्म: $\tau_{12}\tau_{13} = c^2$, जबकि
$\tau_{13}\tau_{12} = c$। (एक प्रविष्टि जाँचने के लिए: $\tau_{12}\tau_{13}$
$1 \xmapsto{\tau_{13}} 3 \xmapsto{\tau_{12}} 3$, $3 \mapsto 1 \mapsto 2$, $2
\mapsto 2 \mapsto 1$ भेजता है: वही $1 \mapsto 3 \mapsto 2 \mapsto 1$ है,
अर्थात् चक्र $c^2 = (1\,3\,2)$।)
\end{solution}

\begin{solution}{exo:b1:structures:4}
$H$: $1 \in H$; $z, w \in H$ के लिए $\abs{zw^{-1}} = \abs z / \abs w = 1$:
कसौटी लागू हो जाती है। $\R_+^*$: वही, बस $\abs{xy^{-1}}$ की जगह धनात्मकता।
सम्मिलन: $\iu \in H$ और $2 \in \R_+^*$, पर $2\iu$ का \omterm{def:b1:complex:field}{मापांक} $2 \neq 1$ है और
वह धनात्मक वास्तविक नहीं है: $2\iu \notin H \cup \R_+^*$, अतः सम्मिलन स्थायी
नहीं है --- \omterm{def:b1:structures:subgroup}{उपसमूह} नहीं (जैसा \cref{exo:b1:structures:6} बताता है, दोनों
उपसमूहों में से कोई दूसरे को नहीं समेटता)।
\end{solution}

\begin{solution}{exo:b1:structures:5}
\omterm{def:b1:structures:morphism}{समाकारिता}: $\eu^{\iu(\theta + \varphi)} = \eu^{\iu\theta}\eu^{\iu\varphi}$
(\cref{thm:b1:complex:funceq})। \omterm{def:b1:structures:morphism}{अष्टि}: $\eu^{\iu\theta} = 1 \iff \theta \in
2\pi\Z$, अतः $\ker f = 2\pi\Z \neq \{0\}$: \omterm{def:b1:logic:inj}{एकैकी} नहीं। प्रतिबिंब: प्रत्येक
एकांक \omterm{def:b1:complex:field}{मापांक} वाली सम्मिश्र संख्या किसी $\theta$ के लिए $\eu^{\iu\theta}$ है
(ध्रुवीय रूप), अतः $\operatorname{im} f = \mathbb{U}$, अर्थात् इकाई वृत्त।
$f$ का $2\pi$ लंबाई के किसी अर्ध-विवृत अंतराल, जैसे $\intco{0}{2\pi}$ या
$\intoc{-\pi}{\pi}$, तक प्रतिबंधन \omterm{def:b1:logic:inj}{एकैकी} है (एक ही प्रतिबिंब वाले दो कोण
$2\pi$ के गुणज से भिन्न होते हैं, और हर वर्ग का केवल एक प्रतिनिधि उस अंतराल
में समाता है); इससे बड़ी लंबाई का कोई अंतराल काम नहीं करता, क्योंकि उसमें
$2\pi$ दूरी पर दो बिंदु आ जाते हैं।
\end{solution}

\begin{solution}{exo:b1:structures:6}
प्रतिच्छेदन: $e \in H \cap K$, और $x, y \in H \cap K$ से $H$ तथा $K$ दोनों
में $xy^{-1}$ मिलता है। सम्मिलन: यदि $H \subseteq K$, तो सम्मिलन $K$ है, जो
\omterm{def:b1:structures:subgroup}{उपसमूह} है (और सममित रूप से भी)। विलोमतः, मान लीजिए कोई भी अंतर्विष्टि सत्य
नहीं है: $h \in H \setminus K$ और $k \in K \setminus H$ चुनिए, और मान लीजिए
$H \cup K$ \omterm{def:b1:structures:subgroup}{उपसमूह} होता; तब $hk \in H \cup K$। यदि $hk \in H$, तो $k =
h^{-1}(hk) \in H$: विरोधाभास। यदि $hk \in K$, तो $h = (hk)k^{-1} \in K$:
विरोधाभास। अतः $H \cup K$ \omterm{def:b1:structures:subgroup}{उपसमूह} नहीं है।
\end{solution}

\begin{solution}{exo:b1:structures:7}
पहले ध्यान दीजिए कि $x^2 = e$ का अर्थ है प्रत्येक $x$ के लिए $x^{-1} = x$।
फिर $x, y \in G$ के लिए:
\[
xy = (xy)^{-1} = y^{-1} x^{-1} = yx ,
\]
जहाँ \cref{prop:b1:structures:rules}~(2) प्रयुक्त हुआ। अतः $G$ \omterm{def:b1:structures:group}{आबेली} है।
\end{solution}

\begin{solution}{exo:b1:structures:8}
$\Z/18\Z$ की इकाइयाँ: $18 = 2 \times 3^2$ से \omterm{cor:b1:arith:bezout}{सहअभाज्य} वर्ग: $\overline 1,
\overline 5, \overline 7, \overline{11}, \overline{13}, \overline{17}$।
$\overline 5$ का प्रतिलोम: $5 \times 11 = 55 = 3\times 18 + 1$, अतः
$\overline 5^{-1} = \overline{11}$।

$\overline 5 x = \overline 7$: $\overline{11}$ से गुणा कीजिए: $x =
\overline{77} = \overline 5$ (क्योंकि $77 = 4\times 18 + 5$)। अद्वितीय हल।

$\overline 6 x = \overline 3$: समीकरण $6x \equiv 3 \pmod{18}$ का अर्थ है $18
\mid 6x - 3$। पर $6x - 3 = 3(2x - 1)$ विषम है, जबकि $18$ सम: कोई सम संख्या
विषम संख्या को विभाजित नहीं कर सकती। कोई हल नहीं।

$\overline 6 x = \overline{12}$: $6x \equiv 12 \pmod{18} \iff x \equiv 2
\pmod 3$: हल $x \in \{\overline 2, \overline 5, \overline 8, \overline{11},
\overline{14}, \overline{17}\}$ --- अर्थात् छह।
\end{solution}

\begin{solution}{exo:b1:structures:9}
$\Z[\sqrt 2]$ में $0$ और $1$ हैं, और वह घटाव तथा गुणनफल के अंतर्गत स्थायी
है:
\[
(a + b\sqrt 2)(c + d\sqrt 2) = (ac + 2bd) + (ad + bc)\sqrt 2 ,
\]
अतः वह $\R$ का उपवलय है (क्रमविनिमेयता, साहचर्यता और वितरणशीलता विरासत में
मिलती हैं)। इकाई: $(1 + \sqrt 2)(-1 + \sqrt 2) = 2 - 1 = 1$, अतः $1 + \sqrt
2$ व्युत्क्रमणीय है और उसका प्रतिलोम $\sqrt 2 - 1 \in \Z[\sqrt 2]$ है। उसकी
घातें $(1 + \sqrt 2)^n$ कठोरतः वर्धमान हैं (आधार $> 1$ है), अतः जोड़ों में
भिन्न हैं, और हर एक इकाई है ($\bigl((1+\sqrt2)^n\bigr)^{-1} = (\sqrt 2 -
1)^n$): अतः इकाइयों का \omterm{def:b1:structures:group}{समूह} अनंत है।
\end{solution}

\begin{solution}{exo:b1:structures:10}
$x + x = (x + x)^2 = x^2 + x^2 + x^2 + x^2 = 4x^2 = 4x$ --- अतः $2x = 4x$,
जिससे $2x = 0$ मिलता है, अर्थात् $x + x = 0$ (हर अवयव अपना ही योगात्मक
प्रतिलोम है)। फिर
\[
x + y = (x+y)^2 = x^2 + xy + yx + y^2 = x + xy + yx + y ,
\]
अतः $xy + yx = 0$, अर्थात् $xy = -yx = yx$ ($-z = z$ का प्रयोग करते हुए)।
अतः $A$ क्रमविनिमेय है।

उदाहरण: $\mathcal{P}(E)$ पर $A + B = (A \cup B) \setminus (A \cap B)$ (सममित
अंतर) और $A \times B = A \cap B$ परिभाषित कीजिए। जाँच लीजिए:
$(\mathcal{P}(E), +)$ \omterm{def:b1:structures:group}{आबेली समूह} है जिसका तत्समक $\emptyset$ है और हर
\omterm{def:b1:logic:sets}{समुच्चय} अपना ही प्रतिलोम; $\cap$ साहचर्य और क्रमविनिमेय है, जिसका तत्समक $E$
है; वितरणशीलता $A \cap (B + C) = (A \cap B) + (A \cap C)$ सत्य है (कोई अवयव
बाएँ पक्ष में तभी है जब वह $A$ में हो और $B, C$ में से ठीक एक में)। और $A
\cap A = A$: हर अवयव वर्गसम है, जैसा अपेक्षित था।
\end{solution}

\begin{solution}{exo:b1:structures:11}
परिकल्पना $n$ और $n+1$ के लिए लिखिए:
\[
(xy)^{n+1} = x^{n+1} y^{n+1}
\quad\text{और}\quad
(xy)^{n+1} = (xy)(xy)^n = xy\,x^n y^n .
\]
बराबर करने पर: $x^{n+1} y^{n+1} = x\,y\,x^n\,y^n$; बाईं ओर $x$ और दाईं ओर
$y^n$ काटिए: $x^n y = y x^n$। एक घात ऊपर वही संगणना ($n+1$ और $n+2$)
$x^{n+1} y = y x^{n+1}$ देती है। फिर
\[
y\,x^{n+1} = x^{n+1} y = x\,(x^n y) = x\,y\,x^n ,
\]
और $y x \cdot x^n = x y \cdot x^n$ की दाईं ओर $x^n$ काटने पर: $yx = xy$। अतः
$G$ \omterm{def:b1:structures:group}{आबेली} है।
\end{solution}

\begin{solution}{exo:b1:structures:12}
\begin{enumerate}
  \item मान लीजिए $f \colon \Z \to \Z$ योगात्मक है और $a = f(1)$। आगमन से $k
  \in \N$ के लिए $f(k) = ka$, और $f(-k) = -f(k) = -ka$: अतः $f$ $a$ से गुणन
  है। विलोमतः प्रत्येक \omterm{def:b1:logic:map}{प्रतिचित्रण} $k \mapsto ak$ एक \omterm{def:b1:structures:morphism}{समाकारिता} है: अर्थात्
  समाकारिताएँ $(\Z,+) \to (\Z,+)$ ठीक किसी नियत पूर्णांक से गुणन हैं।
  \item मान लीजिए $f \colon \Q \to \Z$ एक \omterm{def:b1:structures:morphism}{समाकारिता} है, $x \in \Q$ और $n \in
  \N^*$। तब
        \[
        f(x) = f\Bigl(\underbrace{\tfrac xn + \dots +
        \tfrac xn}_{n}\Bigr) = n\,f\Bigl(\frac xn\Bigr) ,
        \]
        अतः पूर्णांक $f(x)$ प्रत्येक $n \geq 1$ से विभाज्य है। ऐसा एकमात्र
        पूर्णांक $0$ है: $f \equiv 0$।
\end{enumerate}
\end{solution}

\begin{solution}{pb:b1:structures:1}
\textbf{1.} \omterm{def:b1:counting:objects}{क्रमचय} $\intint1n$ का एकैकी आच्छादन है, अर्थात् $n$ वस्तुओं का
$n$-विन्यास: उनकी संख्या $n!$ है (\cref{thm:b1:counting:counts})। $\sigma =
[2,3,1]$, $\tau = [1,3,2]$ के साथ: $\sigma\tau$ $1 \mapsto 1 \mapsto 2$, $2
\mapsto 3 \mapsto 1$, $3 \mapsto 2 \mapsto 3$ भेजता है: $\sigma\tau =
[2,1,3]$; और $\tau\sigma$ $1 \mapsto 2 \mapsto 3$, $2 \mapsto 3 \mapsto 2$,
$3 \mapsto 1 \mapsto 1$ भेजता है: $\tau\sigma = [3,2,1] \neq \sigma\tau$।

\textbf{2.} मान लीजिए $\gamma, \gamma'$ के आश्रय $S, S'$ असंयुक्त हैं। $x
\in S$ के लिए: $\gamma'(x) = x$ और $\gamma(x) \in S$, अतः $\gamma\gamma'(x)
= \gamma(x) = \gamma'\gamma(x)$। $x \in S'$ के लिए सममित रूप से; और दोनों
पक्ष प्रत्येक $x \notin S \cup S'$ को अचल रखते हैं। अतः $\gamma\gamma' =
\gamma'\gamma$।

\textbf{3.} $i \in \intint1n$ के लिए मान $i, \sigma(i), \sigma^2(i), \dots$
किसी \omterm{def:b1:counting:card}{परिमित समुच्चय} में रहते हैं, अतः किसी $a < b$ के लिए $\sigma^a(i) =
\sigma^b(i)$; एकैकीयता $\sigma^{b-a}(i) = i$ देती है: अनुक्रम $i$ पर लौट आता
है। $i$ की \emph{कक्षा} उस \omterm{def:b1:logic:sets}{समुच्चय} $\{i, \sigma(i), \dots,
\sigma^{k-1}(i)\}$ को कहिए, जहाँ $k \geq 1$ ऐसा न्यूनतम है कि $\sigma^k(i) =
i$। एक बिंदु पर मिलने वाली दो कक्षाएँ संपाती हो जाती हैं (दोनों उस बिंदु के
अग्रगामी $\sigma$-प्रतिबिंब हैं), अतः कक्षाएँ $\intint1n$ का \omterm{thm:b1:logic:partition}{विभाजन} करती
हैं; $\sigma$ आकार $k \geq 2$ की प्रत्येक कक्षा पर $k$-चक्र $(i\ \sigma(i)\
\cdots\ \sigma^{k-1}(i))$ की भाँति क्रिया करता है और एकल कक्षाओं को अचल रखता
है। इन असंयुक्त चक्रों का गुणनफल सर्वत्र $\sigma$ से मेल खाता है।
अद्वितीयता: असंयुक्त चक्रों में किसी भी वियोजन में $i$ से होकर जाने वाला
चक्र $(i\ \sigma(i)\ \cdots)$ ही होना चाहिए --- चक्र अपनी प्रेरित क्रिया के
साथ कक्षाएँ होने पर बाध्य हैं।

\textbf{4.} कक्षाओं के अनुसार: $1 \to 4 \to 2 \to 1$, $3 \to 5 \to 3$, $6
\to 7 \to 8 \to 6$:
\[
\sigma = (1\ 4\ 2)(3\ 5)(6\ 7\ 8) .
\]
यदि $\sigma = \gamma_1\cdots\gamma_r$, जहाँ असंयुक्त चक्रों की लंबाइयाँ
$k_1, \dots, k_r$ हैं, तो क्रम-विनिमेयता (प्रश्न 2) $\sigma^m =
\gamma_1^m\cdots\gamma_r^m$ देती है, और चूँकि आश्रय असंयुक्त हैं, $\sigma^m
= \mathrm{id}$ तभी जब प्रत्येक $\gamma_i^m = \mathrm{id}$, अर्थात् तभी जब
सभी $i$ के लिए $k_i \mid m$ ($k$-चक्र की \omterm{def:b1:structures:order}{कोटि} $k$ होती है: $\gamma^m$ $a_1$
को $a_{1 + (m \bmod k)}$ पर भेजता है)। ऐसा लघुतम $m$
$\operatorname{lcm}(k_1, \dots, k_r)$ है। यहाँ: $\operatorname{lcm}(3, 2, 3)
= 6$।

\textbf{5.} दाएँ पक्ष को प्रत्येक बिंदु पर लगाइए, सबसे दाएँ गुणनखंड से आरंभ
करते हुए। $(a_1\ a_2)$ से $a_1 \mapsto a_2$, फिर हर बाद का गुणनखंड $a_2$ को
अचल रखता है: शुद्ध परिणाम $a_1 \mapsto a_2$। $2 \leq i < k$ के लिए: $a_i$ तब
तक अछूता रहता है जब तक $(a_1\ a_i)$ उसे $a_1$ पर नहीं भेजता, और ठीक अगला
गुणनखंड $(a_1\ a_{i+1})$ $a_1$ को $a_{i+1}$ पर भेज देता है, जिसके बाद कोई
उसे नहीं हिलाता: शुद्ध परिणाम $a_i \mapsto a_{i+1}$। अंत में $a_k$ को सबसे
बाएँ वाले के अतिरिक्त सभी गुणनखंड अचल रखते हैं, और वह उसे $a_1$ पर भेज देता
है। यह ठीक वही चक्र है। चूँकि प्रत्येक \omterm{def:b1:counting:objects}{क्रमचय} चक्रों का गुणनफल है (प्रश्न
3), वह व्यत्यासों का गुणनफल भी है। प्रश्न 4 के $\sigma$ के लिए:
\[
\sigma = (1\ 2)(1\ 4)\;(3\ 5)\;(6\ 8)(6\ 7),
\]
पाँच व्यत्यास।

\textbf{6.} $b - a$ पर आगमन। $b = a + 1$ के लिए सर्वसमिका तुच्छ है ($1 =
2\cdot1 - 1$ गुणनखंड)। $b > a + 1$ के लिए सीधे सत्यापित कीजिए कि $(a\ b) =
(a\ \ a{+}1)\,(a{+}1\ \ b)\,(a\ \ a{+}1)$: दायाँ पक्ष $a \mapsto a{+}1
\mapsto b \mapsto b$, $b \mapsto b \mapsto a{+}1 \mapsto a$, $a{+}1 \mapsto
a \mapsto a \mapsto a{+}1$ भेजता है और शेष को अचल रखता है। आगमन से $(a{+}1\
\ b)$ $2(b - a - 1) - 1$ संलग्न व्यत्यासों का पूर्वापर-सम गुणनफल है, अतः
$(a\ b)$ $2(b - a) - 1$ में से एक है: अर्थात् विषम संख्या।

\textbf{7.} $N(\mathrm{id}) = 0$, $\varepsilon = +1$। $(i\ \ i{+}1)$ के लिए
एकमात्र व्युत्क्रमित युग्म $(i, i+1)$ है: $N = 1$, $\varepsilon = -1$। $[2,
3, 1]$ के लिए: व्युत्क्रमित युग्म $(1, 3)$ (मान $2 > 1$) और $(2, 3)$ (मान $3
> 1$) हैं: $N = 2$, $\varepsilon = +1$।

\textbf{8.} $\sigma$ और $\sigma\tau$ की मान-सूचियाँ केवल स्थानों $i$ और $i +
1$ की अदला-बदली से भिन्न हैं। जिस युग्म में $i, i+1$ नहीं है, उसमें कुछ नहीं
बदलता। $k < i$ के लिए दोनों युग्म $(k, i)$ और $(k, i+1)$ अपनी
व्युत्क्रमण-स्थिति की अदला-बदली कर लेते हैं (वही दो मान $\sigma(k)$ से तुलित
होते हैं, बस स्थानों का क्रम उलटा): अतः उनका कुल योगदान अपरिवर्तित रहता है;
$k > i + 1$ के लिए भी वही। शेष बचा एकमात्र युग्म $(i, i+1)$ अपनी स्थिति पलट
देता है। अतः $N(\sigma\tau) = N(\sigma) \pm 1$।

\textbf{9.} मान लीजिए $\tau = (a\ b)$ कोई व्यत्यास है: प्रश्न 6 से वह संलग्न
व्यत्यासों की विषम संख्या का गुणनफल है, अतः दाईं ओर से $\tau$ से गुणा करने
पर $N$ विषम कुल से बदलता है (प्रश्न 8, बार-बार लगाकर): $\varepsilon(\sigma
\tau) = -\varepsilon(\sigma)$। अब यदि $\sigma = \tau_1\cdots \tau_p$
(व्यत्यास), तो उसे तत्समक से $p$ बार दाईं ओर गुणा करके बनाइए:
$\varepsilon(\sigma) = (-1)^p\varepsilon(\mathrm{id}) = (-1)^p$। चूँकि
$\varepsilon(\sigma)$ व्युत्क्रमणों से परिभाषित है --- किसी भी गुणनखंडन से
स्वतंत्र --- अतः $p$ की सम-विषमता $\sigma$ का अपरिवर्त्य है। \omterm{def:b1:structures:morphism}{समाकारिता}: $p$
व्यत्यासों के साथ $\sigma$ और $q$ व्यत्यासों के साथ $\sigma'$ लिखने पर
$\sigma\sigma'$ उनमें से $p + q$ का प्रयोग करता है:
$\varepsilon(\sigma\sigma') = (-1)^{p+q} =
\varepsilon(\sigma)\varepsilon(\sigma')$।

\textbf{10.} $k$-चक्र $k - 1$ व्यत्यासों का गुणनफल है (प्रश्न 5):
$\varepsilon = (-1)^{k-1}$। व्यापक $\sigma$ के लिए, जिसकी कक्षाओं के आकार
$k_1, \dots, k_r$ ($k_i \geq 2$) हैं और साथ $f$ अचल बिंदु हैं, $c(\sigma) =
r + f$ और $n = k_1 + \dots + k_r + f$, अतः
\[
\varepsilon(\sigma) = \prod_{i=1}^r (-1)^{k_i - 1}
= (-1)^{\sum_i k_i - r} = (-1)^{n - f - r} = (-1)^{n -
c(\sigma)} .
\]

\textbf{11.} $\mathfrak A_n = \ker\varepsilon$ एक \omterm{def:b1:structures:morphism}{समाकारिता} की \omterm{def:b1:structures:morphism}{अष्टि} के रूप
में \omterm{def:b1:structures:subgroup}{उपसमूह} है (\cref{def:b1:structures:morphism})। एक व्यत्यास $\tau_0$ नियत
कीजिए ($n \geq 2$ के लिए विद्यमान)। \omterm{def:b1:logic:map}{प्रतिचित्रण} $\sigma \mapsto
\sigma\tau_0$ $\mathfrak S_n$ का एकैकी आच्छादन है (अपना ही प्रतिलोम) जो
$\mathfrak A_n$ और विषम क्रमचयों के \omterm{def:b1:logic:sets}{समुच्चय} की अदला-बदली करता है (प्रश्न 9)।
दोनों \omterm{def:b1:logic:sets}{समुच्चय} $\mathfrak S_n$ का \omterm{thm:b1:logic:partition}{विभाजन} करते हैं और उनके आकार बराबर हैं:
$\abs{\mathfrak A_n} = \frac{n!}2$।

\textbf{12.} $[4, 1, 5, 2, 3, 7, 8, 6]$ के \emph{व्युत्क्रमण}: मान $4$ से:
$1, 2, 3$ पर: तीन; $5$ से: $2, 3$ पर: दो; $7$ से: $6$ पर: एक; $8$ से: $6$
पर: एक। $N = 7$, $\varepsilon = -1$। \emph{चक्र-प्रकार}: $c = 3$ कक्षाएँ, $n
= 8$: $\varepsilon = (-1)^{8-3} = -1$। \emph{व्यत्यास-गणना}: प्रश्न 5 में
पाँच व्यत्यास: $(-1)^5 = -1$। तीनों मेल खाते हैं।

\textbf{13.} $(a\ b)(a\ c)$ (सबसे दाएँ से आरंभ): $a \mapsto c \mapsto c$; $c
\mapsto a \mapsto b$; $b \mapsto b \mapsto a$: अर्थात् $3$-चक्र $(a\ c\ b)$।
और $(a\ c\ b)(a\ c\ d)$: $a \mapsto c \mapsto b$; $b \mapsto b \mapsto a$;
$c \mapsto d \mapsto d$; $d \mapsto a \mapsto c$: अर्थात् $(a\ b)(c\ d)$,
जैसा दावा किया गया था।

\textbf{14.} मान लीजिए $\sigma \in \mathfrak A_n$: प्रश्न 9 से $\sigma =
\tau_1\cdots\tau_{2m}$, जिसमें व्यत्यासों की संख्या सम है। उन्हें क्रमागत
जोड़ों $\tau_{2i-1}\tau_{2i}$ में समूहबद्ध कीजिए: यदि दोनों बराबर हैं, तो
जोड़ा तत्समक है और लुप्त हो जाता है; यदि वे ठीक एक बिंदु साझा करते हैं, तो
प्रश्न 13 की पहली सर्वसमिका जोड़े को एक $3$-चक्र के रूप में लिख देती है; और
यदि वे असंयुक्त हैं, तो दूसरी सर्वसमिका उसे दो $3$-चक्रों के रूप में। अतः
$\sigma$ $3$-चक्रों का गुणनफल है (अथवा तत्समक, जो रिक्त गुणनफल है --- और $n
\geq 3$ के लिए $(1\ 2\ 3)^3$ भी)।

\textbf{15.} $(1\ 2)(3\ 4) = (1\ 3\ 2)(1\ 3\ 4)$ ($a{=}1, b{=}2, c{=}3,
d{=}4$ के साथ प्रश्न 13)। $5$-चक्र के लिए: प्रश्न 5 से $(1\ 2\ 3\ 4\ 5) =
(1\ 5)(1\ 4)(1\ 3)(1\ 2)$, और युग्मन करने पर: $(1\ 5)(1\ 4) = (1\ 4\ 5)$,
$(1\ 3)(1\ 2) = (1\ 2\ 3)$:
\[
(1\ 2\ 3\ 4\ 5) = (1\ 4\ 5)(1\ 2\ 3) .
\]
($3$ पर जाँच: $(1\ 2\ 3)$ $3 \to 1$ भेजता है, फिर $(1\ 4\ 5)$ $1 \to 4$
भेजता है: शुद्ध परिणाम $3 \to 4$, जो सही है।)

\textbf{16.} दोनों पक्षों को किसी स्वेच्छ बिंदु पर लगाइए। $i = \sigma(a_j)$
के लिए: बायाँ पक्ष $\sigma\bigl((a_1\ \dots\ a_k)(a_j)\bigr) =
\sigma(a_{j+1})$ देता है (सूचकांक $k$ के सापेक्ष), और दायाँ पक्ष
$\sigma(a_j)$ के साथ यही करता है। जो $i$ इस रूप का नहीं है, उसके लिए:
$\sigma^{-1}(i)$ आश्रय के बाहर है, अतः बायाँ पक्ष $i$ को अचल रखता है, और
दायाँ पक्ष भी। सर्वत्र बराबर।

\textbf{17.} कोष्ठ $c'$ की पट्टिका को रिक्त कोष्ठ $c$ में सरकाने से कोष्ठों
$c$ और $c'$ की सामग्री की अदला-बदली हो जाती है (पट्टिका $9$, अर्थात् रिक्त
स्थान, $c'$ पर चला जाता है)। यदि पट्टिका $\sigma(i)$ कोष्ठ $i$ में बैठी थी,
तो नई स्थिति $\sigma' = \sigma \circ (c\ c')$ है: वही सामग्री, केवल कोष्ठ
$c, c'$ एक-दूसरे की पुरानी सामग्री पढ़ रहे हैं। प्रश्न 9 से
$\varepsilon(\sigma') = -\varepsilon(\sigma)$।

\textbf{18.} कोई चाल रिक्त कोष्ठ को सटे हुए कोष्ठ पर भेजती है: उसकी पंक्ति
या उसका स्तंभ ठीक $1$ से बदलता है, अतः कोष्ठ $9$ तक की पंक्ति-स्तंभ दूरी $d$
$\pm1$ से बदलती है, और $(-1)^d$ पलट जाता है। चूँकि हर चाल
$\varepsilon(\sigma)$ और $(-1)^{d(\sigma)}$ दोनों को पलटती है, उनका गुणनफल
$I(\sigma)$ हर चाल से अपरिवर्तित रहता है: अर्थात् अपरिवर्त्य।

\textbf{19.} हल की हुई स्थिति में $\varepsilon = +1$, $d = 0$: $I = +1$।
लक्ष्य (पट्टिकाएँ $7, 8$ अदला-बदली, रिक्त स्थान घर पर) कोष्ठों $7$ और $8$ की
सामग्री का व्यत्यास है: $\varepsilon = -1$, $d = 0$: $I = -1$। चूँकि $I$
अपरिवर्त्य है और दोनों मान भिन्न हैं, चालों का कोई अनुक्रम उन्हें नहीं जोड़
सकता।

\textbf{20.} रिक्त स्थान घर पर होने वाली स्थिति कोष्ठों $1, \dots, 8$ में
$8$ पट्टिकाओं का \omterm{def:b1:counting:objects}{क्रमचय} है, अर्थात् $\mathfrak S_8$ का अवयव; उसमें $d = 0$
है, अतः $I = \varepsilon(\sigma)$। पहुँच सकने के लिए $I = +1$ आवश्यक है,
अर्थात् $\sigma \in \mathfrak A_8$; और स्वीकृत विलोम कहता है कि पूरा
$\mathfrak A_8$ पाया जा सकता है। गणना: $\abs{\mathfrak A_8} = \frac{8!}2 =
20\,160$ (प्रश्न 11)।

\textbf{21.} प्रश्न 20 से रिक्त-स्थान-घर पर वाले पहुँच सकने योग्य विन्यास
ठीक $\mathfrak A_8$ बनाते हैं --- विशेष रूप से $\mathfrak S_8$ का एक
\emph{\omterm{def:b1:structures:subgroup}{उपसमूह}}: दो हल हो सकने वाले मिश्रणों का संयोजन, या किसी एक का
प्रतिलोमन, हल हो सकने वाला ही रहता है, जो शुद्ध पहेली-तर्क से स्पष्ट होने से
बहुत दूर है।

\textbf{22.} (क) रिक्त स्थान घर पर रखते हुए पट्टिकाओं का $3$-चक्र:
$\varepsilon = +1$ (प्रश्न 10), $d = 0$, अतः $I = +1$: पहुँचा जा सकता है
(स्वीकृत विलोम से) --- अर्थात् तीन पट्टिकाओं को चक्रित किया जा सकता है। (ख)
पट्टिका $5$ और रिक्त स्थान की अदला-बदली: यह स्थिति कोष्ठों $5$ और $9$ की
सामग्री का व्यत्यास है, अतः $\varepsilon = -1$; रिक्त स्थान केंद्र पर बैठा
है, अपने घर से पंक्ति-स्तंभ दूरी $d = 2$ पर, अतः $(-1)^d = +1$ और $I = -1$:
यहाँ नहीं पहुँचा जा सकता। ``रिक्त स्थान को बीच में खड़ा कर दीजिए'' और शेष
पट्टिकाएँ क्रम में रहें --- यह इतना सरल नहीं है।

\textbf{23.} मान लीजिए $f \colon \mathfrak S_n \to \{\pm1\}$ एक \omterm{def:b1:structures:morphism}{समाकारिता}
है। किन्हीं दो व्यत्यासों $\tau, \tau'$ के लिए प्रश्न 16
$\sigma\tau\sigma^{-1} = \tau'$ वाला $\sigma$ दे देता है (दोनों हिलने वाले
बिंदुओं को दूसरे दोनों पर भेजिए; $n \geq 3$ इसके लिए स्थान की गारंटी देता
है, यद्यपि यहाँ $n = 2$ भी तुच्छ है)। तब $f(\tau') =
f(\sigma)f(\tau)f(\sigma)^{-1} = f(\tau)$, क्योंकि $\{\pm1\}$ \omterm{def:b1:structures:group}{आबेली} है:
अर्थात् $f$ व्यत्यासों पर अचर है। यदि वह अचर $+1$ है, तो व्यत्यासों के सभी
गुणनफलों पर $f = 1$, अर्थात् सर्वत्र (प्रश्न 5)। यदि वह $-1$ है, तो $p$
व्यत्यासों के गुणनफल पर $f(\sigma) = (-1)^p = \varepsilon(\sigma)$। अतः $f
\in \{1, \varepsilon\}$।

\textbf{24.} (क) $\varepsilon$ के \omterm{def:b1:structures:morphism}{समाकारिता} होने और \omterm{def:b1:structures:morphism}{अष्टि} की मशीनरी ने
$\mathfrak A_n$ को उसकी \omterm{def:b1:structures:subgroup}{उपसमूह} संरचना और उसका आकार दिया, और
\cref{prop:b1:structures:kernel} की शैली का तर्क प्रश्न 11 तथा 21 में चलता
है। (ख) गणना: $\abs{\mathfrak S_n} = n!$, प्रश्न 11 का आधा करने वाला तर्क,
और प्रश्न 20 की गणना $20\,160$ --- ये सब \cref{ch:b1:counting} काम पर हैं।
(ग) प्रश्न 8--9 एक सच्ची सुपरिभाषितता की समस्या सुलझाते हैं --- ``व्यत्यासों
की संख्या की सम-विषमता'' यह पूर्वमान लेती है कि यह सम-विषमता गुणनखंडन पर
निर्भर नहीं करती, ठीक वैसे ही जैसे $\Z/n\Z$ की संक्रियाओं को
\cref{def:b1:structures:zn} में प्रतिनिधि-निरपेक्षता चाहिए थी।

\textbf{25.} संकेतक एक ही $\{\pm1\}$-मान वाली संगणना है, जिसकी सुपरिभाषितता
एक बार सिद्ध हो जाती है, और वह एक साथ तीन काम करता है: भीतर से वह $\mathfrak
S_n$ को आधा काट देता है और $3$-चक्र जनकों वाले $\mathfrak A_n$ को अलग कर
देता है; बाहर से वह एक ही पंक्ति में ऐसा प्रश्न तय कर देता है (``क्या इन दो
पट्टिकाओं की अदला-बदली संभव है?'') जिसे भोली खोज कभी नहीं सुलझा सकती थी,
क्योंकि विफल चाल-अनुक्रमों की कोई परिमित सूची असंभवता सिद्ध नहीं करती; और
संरचनात्मक रूप से वह \cref{ch:b1:det} के सूत्र $\det A = \sum_\sigma
\varepsilon(\sigma)\, a_{1\sigma(1)}\cdots a_{n\sigma(n)}$ के भीतर बैठा
एकांतरित-चिह्न इंजन है। यह प्रतिरूप --- ऐसी राशि खोजिए जो हर प्रारंभिक चाल
से संरक्षित रहती हो, फिर उसे आरंभ पर और लक्ष्य पर संगणित कीजिए --- ``क्या यह
संभव है?'' जैसे प्रश्नों के विरुद्ध गणितज्ञ का मानक अस्त्र है, और जब भी कोई
\omterm{def:b1:structures:group}{समूह} अवस्थाओं के \omterm{def:b1:logic:sets}{समुच्चय} पर क्रिया करेगा, वह लौट आएगा।
\end{solution}
